\documentclass[11pt,a4paper]{article}

\usepackage[T1]{fontenc}
\usepackage[utf8]{inputenc}
\usepackage{lmodern}
\usepackage{microtype}
\usepackage{geometry}
\geometry{margin=1in}
\usepackage{amsmath,amssymb,amsthm,mathtools}
\usepackage{enumitem}
\usepackage{xcolor}
\usepackage{hyperref}
\hypersetup{
  colorlinks=true,
  linkcolor=blue!55!black,
  citecolor=blue!55!black,
  urlcolor=blue!55!black,
  pdftitle={Ext Groups and Ext Sheaves: A Seminar Report on Hartshorne III.6},
  pdfauthor={Chang Wang},
  pdfsubject={Seminar report},
  pdfkeywords={Ext, sheaf Ext, Hartshorne, algebraic geometry}
}

\numberwithin{equation}{section}

\theoremstyle{plain}
\newtheorem{theorem}{Theorem}[section]
\newtheorem{proposition}[theorem]{Proposition}
\newtheorem{lemma}[theorem]{Lemma}
\newtheorem{corollary}[theorem]{Corollary}

\theoremstyle{definition}
\newtheorem{definition}[theorem]{Definition}
\newtheorem{example}[theorem]{Example}

\theoremstyle{remark}
\newtheorem{remark}[theorem]{Remark}

\newcommand{\OO}{\mathcal O}
\newcommand{\F}{\mathcal F}
\newcommand{\G}{\mathcal G}
\newcommand{\HH}{\mathcal H}
\newcommand{\E}{\mathcal E}
\newcommand{\R}{\mathcal R}
\newcommand{\LL}{\mathcal L}
\newcommand{\I}{\mathcal I}
\newcommand{\A}{\mathcal A}
\newcommand{\C}{\mathcal C}
\DeclareMathOperator{\Ext}{Ext}
\DeclareMathOperator{\Hom}{Hom}
\DeclareMathOperator{\coker}{coker}
\newcommand{\sExt}{\mathcal{E}\!xt}
\newcommand{\sHom}{\mathcal{H}\!om}
\newcommand{\Mod}{\operatorname{Mod}}

\title{Ext Groups and Ext Sheaves\\
\large A Seminar Report on Hartshorne III, Section 6}
\author{Chang Wang}
\date{July 2026}

\begin{document}
\maketitle

\begin{abstract}
This report explains the distinction between global Ext groups and sheaf Ext in the category of sheaves of modules on a ringed space, with special attention to coherent sheaves on noetherian schemes.  The main reference is Hartshorne, \emph{Algebraic Geometry}, Chapter III, Section 6.  The central objects are the global groups
\[
  \Ext_X^i(\F,\G)
\]
and the local sheaves
\[
  \sExt_X^i(\F,\G).
\]
They agree in degree zero after taking global sections, but they differ in higher degree because global Ext also contains sheaf cohomology.  The final part proves Hartshorne III, Proposition 6.9: if \(X\) is projective over a noetherian ring and \(\F,\G\) are coherent, then for each fixed \(i\) and all sufficiently large \(n\), the natural map
\[
  \Ext_X^i(\F,\G(n))
  \longrightarrow
  \Gamma\bigl(X,\sExt_X^i(\F,\G(n))\bigr)
\]
is an isomorphism.  This is explained both through the local-to-global Ext spectral sequence and through Hartshorne's dimension-shifting proof.
\end{abstract}

\section{Introduction}

Let \(X\) be a ringed space.  The purpose of this report is to compare two derived functors associated to two sheaves of \(\OO_X\)-modules \(\F\) and \(\G\).  The first is the global Ext group \(\Ext_X^i(\F,\G)\), an abelian group.  The second is the sheaf Ext \(\sExt_X^i(\F,\G)\), a sheaf of \(\OO_X\)-modules.  The notation is deliberately similar, but the two objects contain different information.

The basic degree-zero identity is
\begin{equation}\label{eq:hom-global-sections}
  \Hom_X(\F,\G)
  =
  \Gamma\bigl(X,\sHom_X(\F,\G)\bigr).
\end{equation}
It is natural to ask whether the same identity remains true after deriving both sides.  In general it does not.  The reason is that deriving the functor \(\Hom_X(\F,-)\) mixes the local sheaf Hom construction with global sections, and global sections are not exact.

The simplest example is obtained by taking \(\F=\OO_X\).  Since
\[
  \Hom_X(\OO_X,\G)=\Gamma(X,\G),
\]
right derived functors give
\[
  \Ext_X^i(\OO_X,\G)=H^i(X,\G).
\]
On the other hand, \(\sHom_X(\OO_X,-)\) is the identity functor, and hence
\[
  \sExt_X^i(\OO_X,\G)=0 \qquad (i>0).
\]
Thus global Ext sees sheaf cohomology, while sheaf Ext is local.

The precise relationship is governed by the local-to-global Ext spectral sequence
\begin{equation}\label{eq:local-to-global-intro}
  E_2^{p,q}=H^p\bigl(X,\sExt_X^q(\F,\G)\bigr)
  \Longrightarrow
  \Ext_X^{p+q}(\F,\G).
\end{equation}
Hartshorne III, Proposition 6.9 says that for coherent sheaves on a projective noetherian scheme, the higher cohomology terms in this spectral sequence disappear after sufficiently positive twisting.  Therefore, in high twist, global Ext becomes the global sections of sheaf Ext.

\section{Definitions and first examples}

Throughout this report, \(\Mod(X)\) denotes the category of \(\OO_X\)-modules.  If \(X\) is equipped with an invertible sheaf \(\OO_X(1)\), then
\[
  \G(n)=\G\otimes_{\OO_X}\OO_X(n),
  \qquad
  \OO_X(n)=\OO_X(1)^{\otimes n}.
\]
For negative \(n\), this has the usual meaning using \(\OO_X(-1)=\OO_X(1)^\vee\).  A locally free sheaf is always assumed to have finite rank.

\begin{definition}[Global Ext]
Fix \(\F\in \Mod(X)\).  The functor
\[
  \Hom_X(\F,-):\Mod(X)\longrightarrow \mathbf{Ab}
\]
is left exact.  Since \(\Mod(X)\) has enough injective objects, its right derived functors are defined.  If \(\G\to \I^\bullet\) is an injective resolution, then
\[
  \Ext_X^i(\F,\G)=H^i\bigl(\Hom_X(\F,\I^\bullet)\bigr).
\]
Equivalently, \(\Ext_X^i(\F,\G)=R^i\Hom_X(\F,-)(\G)\).
\end{definition}

\begin{definition}[Sheaf Ext]
The sheaf Hom functor
\[
  \sHom_X(\F,-):\Mod(X)\longrightarrow \Mod(X)
\]
is also left exact.  Its right derived functors are denoted by \(\sExt_X^i(\F,-)\).  Thus, for an injective resolution \(\G\to\I^\bullet\),
\[
  \sExt_X^i(\F,\G)
  =
  H^i\bigl(\sHom_X(\F,\I^\bullet)\bigr),
\]
where \(H^i\) means the \(i\)-th cohomology sheaf.
\end{definition}

For \(i=0\), the two constructions recover the usual Hom objects:
\[
  \Ext_X^0(\F,\G)=\Hom_X(\F,\G),
  \qquad
  \sExt_X^0(\F,\G)=\sHom_X(\F,\G).
\]
The identity \eqref{eq:hom-global-sections} therefore says exactly that global Ext and sheaf Ext agree in degree zero after taking global sections.

\begin{remark}[The comparison map]
For every \(i\) there is a natural edge map
\begin{equation}\label{eq:comparison-map}
  \eta_{\F,\G}^i:
  \Ext_X^i(\F,\G)
  \longrightarrow
  \Gamma\bigl(X,\sExt_X^i(\F,\G)\bigr).
\end{equation}
It can be described by choosing an injective resolution \(\G\to\I^\bullet\).  A class in \(\Ext_X^i(\F,\G)\) is represented by a global cocycle in \(\Gamma(X,\sHom_X(\F,\I^i))\).  Taking its local cohomology class gives a global section of the cohomology sheaf \(\sExt_X^i(\F,\G)\).  The map \eqref{eq:comparison-map} is usually not an isomorphism.
\end{remark}

\begin{proposition}[Ext from the structure sheaf]\label{prop:structure-sheaf}
For every \(\OO_X\)-module \(\G\),
\[
  \sExt_X^0(\OO_X,\G)\cong\G,
  \qquad
  \sExt_X^i(\OO_X,\G)=0\quad (i>0),
\]
and
\[
  \Ext_X^i(\OO_X,\G)\cong H^i(X,\G).
\]
\end{proposition}

\begin{proof}
The functor \(\sHom_X(\OO_X,-)\) is naturally isomorphic to the identity functor on \(\Mod(X)\), so it is exact.  This proves the statement for sheaf Ext.  For global Ext, the functor \(\Hom_X(\OO_X,-)\) is the global section functor \(\Gamma(X,-)\).  Its right derived functors are sheaf cohomology.
\end{proof}

\begin{example}\label{ex:curve}
Let \(X\) be a smooth projective curve of positive genus over a field.  Then \(H^1(X,\OO_X)\ne0\).  Proposition \ref{prop:structure-sheaf} gives
\[
  \Ext_X^1(\OO_X,\OO_X)=H^1(X,\OO_X)\ne0,
  \qquad
  \Gamma\bigl(X,\sExt_X^1(\OO_X,\OO_X)\bigr)=0.
\]
Thus global Ext cannot in general be recovered from the global sections of sheaf Ext.
\end{example}

\section{Basic properties of sheaf Ext}

The following results correspond to Hartshorne III, Propositions 6.2--6.8.  They explain in what sense sheaf Ext is local, while global Ext also contains cohomological information on \(X\).

\begin{proposition}[Restriction to open subsets]\label{prop:restriction}
Let \(U\subseteq X\) be open.  Then there is a natural isomorphism
\[
  \sExt_X^i(\F,\G)|_U
  \cong
  \sExt_U^i(\F|_U,\G|_U).
\]
\end{proposition}

\begin{proof}
Let \(j:U\hookrightarrow X\) be the inclusion.  If \(\I\) is an injective \(\OO_X\)-module, then \(\I|_U\) is injective as an \(\OO_U\)-module.  Indeed, extension by zero \(j_!\) is exact and is left adjoint to restriction, so
\[
  \Hom_U(\A,\I|_U)
  \cong
  \Hom_X(j_!\A,\I)
\]
is an exact functor of \(\A\).  Hence the restriction of an injective resolution \(\G\to\I^\bullet\) is an injective resolution \(\G|_U\to\I^\bullet|_U\).  Since
\[
  \sHom_X(\F,\I^j)|_U
  \cong
  \sHom_U(\F|_U,\I^j|_U)
\]
for each \(j\), taking cohomology sheaves gives the result.
\end{proof}

\begin{proposition}[Long exact sequence in the first variable]\label{prop:les-first-variable}
Let
\[
  0\longrightarrow \F'\longrightarrow \F\longrightarrow \F''\longrightarrow0
\]
be a short exact sequence of \(\OO_X\)-modules.  Then there are natural long exact sequences
\[
\begin{aligned}
0&\longrightarrow \Hom_X(\F'',\G)
 \longrightarrow \Hom_X(\F,\G)
 \longrightarrow \Hom_X(\F',\G) \\
&\longrightarrow \Ext_X^1(\F'',\G)
 \longrightarrow \Ext_X^1(\F,\G)
 \longrightarrow \Ext_X^1(\F',\G)
 \longrightarrow\cdots,
\end{aligned}
\]
and
\[
\begin{aligned}
0&\longrightarrow \sHom_X(\F'',\G)
 \longrightarrow \sHom_X(\F,\G)
 \longrightarrow \sHom_X(\F',\G) \\
&\longrightarrow \sExt_X^1(\F'',\G)
 \longrightarrow \sExt_X^1(\F,\G)
 \longrightarrow \sExt_X^1(\F',\G)
 \longrightarrow\cdots .
\end{aligned}
\]
\end{proposition}

\begin{proof}
Choose an injective resolution \(\G\to\I^\bullet\).  For an injective sheaf \(\I\), the functor \(\Hom_X(-,\I)\) is exact in the first variable.  The sheaf version is also exact: by Proposition \ref{prop:restriction}, \(\I|_U\) is injective for every open \(U\), so \(\Hom_U(-,\I|_U)\) is exact on sections over \(U\).  Applying \(\Hom_X(-,\I^\bullet)\) and \(\sHom_X(-,\I^\bullet)\) therefore gives short exact sequences of complexes.  The associated long exact cohomology sequences are the desired ones.  The arrows reverse because Hom is contravariant in the first variable.
\end{proof}

\begin{proposition}[Computation by locally free resolutions]\label{prop:locally-free-resolution}
Suppose that \(\F\) has a resolution by locally free sheaves of finite rank,
\[
  \cdots\longrightarrow \LL_2\longrightarrow \LL_1
  \longrightarrow \LL_0\longrightarrow \F\longrightarrow0.
\]
Then
\[
  \sExt_X^i(\F,\G)
  \cong
  h^i\bigl(\sHom_X(\LL_\bullet,\G)\bigr),
\]
where the right side denotes the cohomology sheaf of the complex
\[
  0\longrightarrow \sHom_X(\LL_0,\G)
  \longrightarrow \sHom_X(\LL_1,\G)
  \longrightarrow \sHom_X(\LL_2,\G)
  \longrightarrow\cdots .
\]
\end{proposition}

\begin{proof}
If \(\LL\) is locally free of finite rank, then
\[
  \sHom_X(\LL,-)\cong \LL^\vee\otimes_{\OO_X}-.
\]
This functor is exact.  Thus a locally free resolution of \(\F\) may be used in the first variable to compute the right derived sheaves of \(\sHom_X(\F,-)\).  Equivalently, after restricting to an open set on which the \(\LL_j\) are free, this is the usual computation of module Ext by a projective resolution.
\end{proof}

\begin{remark}\label{rem:global-not-computed}
Proposition \ref{prop:locally-free-resolution} computes sheaf Ext, not global Ext.  For instance, \(\OO_X\) has a locally free resolution of length zero, but Proposition \ref{prop:structure-sheaf} gives \(\Ext_X^i(\OO_X,\G)=H^i(X,\G)\), which can be nonzero for \(i>0\).  The missing ingredient is the non-exactness of the global section functor.
\end{remark}

\begin{lemma}[Tensoring injectives with locally free sheaves]\label{lem:injective-tensor}
Let \(\LL\) be locally free of finite rank.  If \(\I\) is an injective \(\OO_X\)-module, then \(\LL\otimes\I\) is injective.
\end{lemma}

\begin{proof}
For every \(\OO_X\)-module \(\A\), there is a natural adjunction
\[
  \Hom_X(\A,\LL\otimes\I)
  \cong
  \Hom_X(\A\otimes\LL^\vee,\I).
\]
The functor \(-\otimes\LL^\vee\) is exact because \(\LL^\vee\) is locally free.  Since \(\I\) is injective, the right side is exact as a functor of \(\A\).  Hence \(\LL\otimes\I\) is injective.
\end{proof}

\begin{proposition}[Moving a locally free factor]\label{prop:moving-locally-free-factor}
Let \(\LL\) be locally free of finite rank.  There are natural isomorphisms
\[
  \Ext_X^i(\F\otimes\LL,\G)
  \cong
  \Ext_X^i(\F,\LL^\vee\otimes\G)
\]
and
\[
  \sExt_X^i(\F\otimes\LL,\G)
  \cong
  \sExt_X^i(\F,\LL^\vee\otimes\G).
\]
\end{proposition}

\begin{proof}
For \(i=0\), this is the usual Hom-tensor adjunction.  For higher \(i\), choose an injective resolution \(\G\to\I^\bullet\).  By Lemma \ref{lem:injective-tensor}, \(\LL^\vee\otimes\I^\bullet\) is an injective resolution of \(\LL^\vee\otimes\G\).  Applying the degree-zero adjunction term by term gives isomorphic complexes, both globally and sheaf-theoretically.  Taking cohomology gives the result.
\end{proof}

\begin{corollary}[Locally free first variable]\label{cor:locally-free-first-variable}
If \(\E\) is locally free of finite rank, then
\[
  \sExt_X^i(\E,\G)=0\qquad(i>0)
\]
and
\[
  \Ext_X^i(\E,\G)
  \cong
  H^i(X,\E^\vee\otimes\G).
\]
\end{corollary}

\begin{proof}
For sheaf Ext, the functor \(\sHom_X(\E,-)\cong \E^\vee\otimes -\) is exact.  For global Ext, Proposition \ref{prop:moving-locally-free-factor} gives
\[
  \Ext_X^i(\E,\G)
  \cong
  \Ext_X^i(\OO_X,\E^\vee\otimes\G),
\]
and Proposition \ref{prop:structure-sheaf} identifies the latter group with \(H^i(X,\E^\vee\otimes\G)\).
\end{proof}

\begin{proposition}[Stalks of sheaf Ext]\label{prop:stalks}
Let \(X\) be a noetherian scheme, let \(\F\) be coherent, and let \(\G\) be an \(\OO_X\)-module.  For every point \(x\in X\), there is a natural isomorphism
\[
  \sExt_X^i(\F,\G)_x
  \cong
  \Ext_{\OO_{X,x}}^i(\F_x,\G_x).
\]
\end{proposition}

\begin{proof}
By Proposition \ref{prop:restriction}, the question is local near \(x\).  We may work on an affine open neighbourhood of \(x\).  Since \(X\) is noetherian and \(\F\) is coherent, \(\F\) has, locally near \(x\), a resolution by finite rank free sheaves.  Proposition \ref{prop:locally-free-resolution} gives
\[
  \sExt_X^i(\F,\G)
  \cong
  h^i\bigl(\sHom_X(\LL_\bullet,\G)\bigr)
\]
near \(x\).  Taking stalks is exact, and for each \(j\),
\[
  \sHom_X(\LL_j,\G)_x
  \cong
  \Hom_{\OO_{X,x}}((\LL_j)_x,\G_x).
\]
The complex \((\LL_\bullet)_x\) is a free resolution of \(\F_x\) over \(\OO_{X,x}\).  Hence the cohomology of the stalk complex is \(\Ext_{\OO_{X,x}}^i(\F_x,\G_x)\).
\end{proof}

\section{The local-to-global spectral sequence}

\begin{theorem}[Local-to-global Ext spectral sequence]\label{thm:local-to-global}
For sheaves \(\F\) and \(\G\) of \(\OO_X\)-modules, there is a spectral sequence
\[
  E_2^{p,q}=H^p\bigl(X,\sExt_X^q(\F,\G)\bigr)
  \Longrightarrow
  \Ext_X^{p+q}(\F,\G).
\]
The edge map in total degree \(i\) is the comparison map
\[
  \Ext_X^i(\F,\G)
  \longrightarrow
  \Gamma\bigl(X,\sExt_X^i(\F,\G)\bigr).
\]
\end{theorem}

\begin{proof}
This is the Grothendieck spectral sequence for the composition
\[
  \Mod(X)
  \xrightarrow{\sHom_X(\F,-)}
  \Mod(X)
  \xrightarrow{\Gamma(X,-)}
  \mathbf{Ab},
\]
since \(\Hom_X(\F,-)=\Gamma(X,\sHom_X(\F,-))\).  The statement about the edge map follows from the construction of the spectral sequence applied to an injective resolution of \(\G\).
\end{proof}

\begin{remark}
For total degree \(i\), the term \(H^0(X,\sExt_X^i(\F,\G))\) is only one part of the spectral sequence.  Other terms such as
\[
  H^1(X,\sExt_X^{i-1}(\F,\G)),\quad
  H^2(X,\sExt_X^{i-2}(\F,\G)),\quad \ldots
\]
may also contribute to \(\Ext_X^i(\F,\G)\).  This is the conceptual reason why the comparison map \eqref{eq:comparison-map} is not usually an isomorphism.
\end{remark}

\section{Coherent sheaves and twists}

From now on assume that \(X\) is projective over a noetherian ring \(A\), equipped with a very ample invertible sheaf \(\OO_X(1)\), and that \(\F\) and \(\G\) are coherent.

\begin{theorem}[Serre vanishing]\label{thm:serre-vanishing}
Let \(\HH\) be a coherent sheaf on \(X\).  Then there exists an integer \(n_0\) such that
\[
  H^p(X,\HH(n))=0
  \qquad
  \text{for all }p>0\text{ and all }n\ge n_0.
\]
\end{theorem}

\begin{lemma}[Exactness after high twisting]\label{lem:high-twist-exactness}
Let
\[
  0\longrightarrow \HH_0\longrightarrow \HH_1\longrightarrow\cdots
  \longrightarrow \HH_r\longrightarrow0
\]
be a finite exact sequence of coherent sheaves on \(X\).  Then, for all \(n\gg0\), the sequence
\[
  0\longrightarrow \Gamma(X,\HH_0(n))
  \longrightarrow \Gamma(X,\HH_1(n))
  \longrightarrow\cdots
  \longrightarrow \Gamma(X,\HH_r(n))
  \longrightarrow0
\]
is exact.
\end{lemma}

\begin{proof}
Break the given finite exact sequence into short exact sequences using successive kernels.  These kernels are coherent because \(X\) is noetherian.  After twisting by \(n\), the only obstruction to exactness on global sections is the first cohomology of one of these coherent kernels.  By Serre vanishing, all such \(H^1\)'s vanish for \(n\gg0\).
\end{proof}

\begin{lemma}[Coherence and twisting of sheaf Ext]\label{lem:twist-ext}
Let \(X\) be a noetherian scheme, let \(\F\) and \(\G\) be coherent sheaves, and let \(\LL\) be an invertible sheaf.  Then \(\sExt_X^i(\F,\G)\) is coherent for every \(i\), and there is a natural isomorphism
\[
  \sExt_X^i(\F,\G\otimes\LL)
  \cong
  \sExt_X^i(\F,\G)\otimes\LL.
\]
In particular,
\[
  \sExt_X^i(\F,\G(n))
  \cong
  \sExt_X^i(\F,\G)(n).
\]
\end{lemma}

\begin{proof}
The assertion is local on \(X\).  Near any point, \(\F\) has a resolution by finite free sheaves.  Applying \(\sHom_X(-,\G)\) gives a complex of coherent sheaves, and the cohomology sheaves are coherent because \(X\) is noetherian.  On an open set where \(\LL\) is trivial, the twisting statement is immediate.  Equivalently, using a local finite free resolution \(\LL_\bullet\) of \(\F\), one has
\[
  \sHom_X(\LL_j,\G\otimes\LL)
  \cong
  \sHom_X(\LL_j,\G)\otimes\LL
\]
for every \(j\).  Since tensoring with an invertible sheaf is exact, this identification passes to cohomology sheaves.
\end{proof}

\section{Hartshorne III, Proposition 6.9}

\begin{theorem}[Hartshorne III, Proposition 6.9]\label{thm:hartshorne-6-9}
Let \(X\) be a projective scheme over a noetherian ring \(A\), let \(\OO_X(1)\) be very ample, and let \(\F\) and \(\G\) be coherent sheaves on \(X\).  For each fixed integer \(i\ge0\), the natural map
\[
  \eta_{\F,\G(n)}^i:
  \Ext_X^i(\F,\G(n))
  \longrightarrow
  \Gamma\bigl(X,\sExt_X^i(\F,\G(n))\bigr)
\]
is an isomorphism for all \(n\gg0\).
\end{theorem}

\begin{proof}[Proof using the local-to-global spectral sequence]
The local-to-global Ext spectral sequence gives
\[
  E_2^{p,q}(n)
  =H^p\bigl(X,\sExt_X^q(\F,\G(n))\bigr)
  \Longrightarrow
  \Ext_X^{p+q}(\F,\G(n)).
\]
By Lemma \ref{lem:twist-ext},
\[
  \sExt_X^q(\F,\G(n))
  \cong
  \sExt_X^q(\F,\G)(n).
\]
Fix the total degree \(i\).  Only the finitely many sheaves \(\sExt_X^q(\F,\G)\) with \(0\le q\le i\) can contribute to total degree \(i\), and each of them is coherent.  By Serre vanishing, after choosing \(n\) sufficiently large for these finitely many sheaves,
\[
  H^p\bigl(X,\sExt_X^q(\F,\G(n))\bigr)=0
  \qquad(p>0,\ 0\le q\le i).
\]
Thus in total degree \(i\) all terms with \(p>0\) vanish.  The term \(E_2^{0,i}\) also cannot be affected by any differential: there is no incoming differential from negative \(p\), and all possible outgoing differentials land in groups with \(p>0\) and \(q\le i\), which vanish for \(n\gg0\).  Hence
\[
  E_\infty^{0,i}=E_2^{0,i}
  =H^0\bigl(X,\sExt_X^i(\F,\G(n))\bigr),
\]
and all other graded pieces of the filtration of \(\Ext_X^i(\F,\G(n))\) vanish.  Therefore the edge map identifies
\[
  \Ext_X^i(\F,\G(n))
  \cong
  H^0\bigl(X,\sExt_X^i(\F,\G(n))\bigr)
\]
for \(n\gg0\).  This edge map is precisely \(\eta_{\F,\G(n)}^i\).
\end{proof}

\begin{proof}[Hartshorne's dimension-shifting proof]
The statement is clear for \(i=0\), because
\[
  \Ext_X^0(\F,\G(n))=
  \Hom_X(\F,\G(n))
  =
  \Gamma\bigl(X,\sHom_X(\F,\G(n))\bigr).
\]
It is also clear for locally free first variable in positive degree after high twisting.  Indeed, if \(\E\) is locally free, then Corollary \ref{cor:locally-free-first-variable} gives
\[
  \Ext_X^i(\E,\G(n))
  \cong
  H^i(X,\E^\vee\otimes\G(n)),
\]
which vanishes for \(i>0\) and \(n\gg0\) by Serre vanishing; while \(\sExt_X^i(\E,\G(n))=0\) for \(i>0\).

Now let \(\F\) be coherent.  Since \(X\) is projective over a noetherian ring, \(\F(m)\) is generated by finitely many global sections for \(m\gg0\).  Hence there is a surjection
\[
  \E=\OO_X(-m)^{\oplus N}\twoheadrightarrow \F
\]
from a locally free sheaf of finite rank.  Let \(\R\) be the kernel:
\[
  0\longrightarrow \R\longrightarrow \E\longrightarrow \F\longrightarrow0.
\]
The sheaf \(\R\) is coherent.

First consider \(i=1\).  The long exact sequence for global Ext gives
\[
\begin{aligned}
0&\longrightarrow \Hom_X(\F,\G(n))
 \longrightarrow \Hom_X(\E,\G(n))
 \longrightarrow \Hom_X(\R,\G(n)) \\
&\longrightarrow \Ext_X^1(\F,\G(n))
 \longrightarrow \Ext_X^1(\E,\G(n)).
\end{aligned}
\]
For \(n\gg0\), the last term vanishes, so
\[
  \Ext_X^1(\F,\G(n))
  \cong
  \coker\bigl(\Hom_X(\E,\G(n))
  \to \Hom_X(\R,\G(n))\bigr).
\]
The corresponding sheaf exact sequence is
\[
\begin{aligned}
0&\longrightarrow \sHom_X(\F,\G(n))
 \longrightarrow \sHom_X(\E,\G(n))
 \longrightarrow \sHom_X(\R,\G(n)) \\
&\longrightarrow \sExt_X^1(\F,\G(n))
 \longrightarrow0,
\end{aligned}
\]
because \(\sExt_X^1(\E,\G(n))=0\).  All sheaves appearing here are twists of fixed coherent sheaves.  By Lemma \ref{lem:high-twist-exactness}, applying \(\Gamma(X,-)\) to this finite exact sequence is exact for \(n\gg0\).  Hence
\[
\begin{aligned}
\Gamma\bigl(X,\sExt_X^1(\F,\G(n))\bigr)
&\cong
\coker\bigl(
\Gamma(X,\sHom_X(\E,\G(n)))
\to
\Gamma(X,\sHom_X(\R,\G(n)))
\bigr) \\
&\cong
\coker\bigl(
\Hom_X(\E,\G(n))
\to
\Hom_X(\R,\G(n))
\bigr) \\
&\cong
\Ext_X^1(\F,\G(n)).
\end{aligned}
\]
This proves the theorem in degree one.

For \(i\ge2\), use induction on \(i\).  The long exact sequence for global Ext contains
\[
  \Ext_X^{i-1}(\E,\G(n))
  \longrightarrow
  \Ext_X^{i-1}(\R,\G(n))
  \longrightarrow
  \Ext_X^i(\F,\G(n))
  \longrightarrow
  \Ext_X^i(\E,\G(n)).
\]
For \(n\gg0\), the two outer terms vanish, and therefore
\[
  \Ext_X^i(\F,\G(n))
  \cong
  \Ext_X^{i-1}(\R,\G(n)).
\]
The sheaf long exact sequence gives, because \(\E\) is locally free,
\[
  \sExt_X^i(\F,\G(n))
  \cong
  \sExt_X^{i-1}(\R,\G(n)).
\]
Taking global sections and applying the induction hypothesis to the coherent sheaf \(\R\) in degree \(i-1\), we obtain
\[
\begin{aligned}
\Ext_X^i(\F,\G(n))
&\cong
\Ext_X^{i-1}(\R,\G(n)) \\
&\cong
\Gamma\bigl(X,\sExt_X^{i-1}(\R,\G(n))\bigr) \\
&\cong
\Gamma\bigl(X,\sExt_X^i(\F,\G(n))\bigr)
\end{aligned}
\]
for \(n\gg0\).  These isomorphisms are compatible with the natural comparison maps, so this proves the theorem.
\end{proof}

\section{Conclusion}

The main point is that \(\sExt_X^i(\F,\G)\) is a local invariant, while \(\Ext_X^i(\F,\G)\) is a global invariant.  The local nature of sheaf Ext is expressed by restriction to open subsets and by the stalk formula
\[
  \sExt_X^i(\F,\G)_x
  \cong
  \Ext_{\OO_{X,x}}^i(\F_x,\G_x)
\]
for coherent \(\F\) on a noetherian scheme.  By contrast, global Ext includes sheaf cohomology: already \(\Ext_X^i(\OO_X,\G)=H^i(X,\G)\).

The local-to-global spectral sequence explains the relation between the two objects.  In general, the terms
\[
  H^p\bigl(X,\sExt_X^q(\F,\G)\bigr)
  \qquad(p+q=i)
\]
all may contribute to \(\Ext_X^i(\F,\G)\).  Hartshorne III, Proposition 6.9 says that on a projective noetherian scheme, after replacing \(\G\) by \(\G(n)\) with \(n\gg0\), Serre vanishing kills all terms with \(p>0\) in each fixed total degree.  Thus the comparison map becomes an isomorphism:
\[
  \Ext_X^i(\F,\G(n))
  \cong
  \Gamma\bigl(X,\sExt_X^i(\F,\G(n))\bigr)
  \qquad(n\gg0).
\]
This result is a useful bridge between the local algebra of Ext sheaves and the global geometry of projective schemes, and it prepares the ground for later applications such as duality theory.

\begin{thebibliography}{9}
\bibitem{Hartshorne}
R. Hartshorne, \emph{Algebraic Geometry}, Graduate Texts in Mathematics 52, Springer, 1977.  See Chapter III, Section 6.
\end{thebibliography}

\end{document}
